% !TEX TS-program = pdflatex
% !TEX encoding = UTF-8 Unicode

% This is a simple template for a LaTeX document using the "article" class.
% See "book", "report", "letter" for other types of document.

\documentclass[10pt]{ltxdoc} % use larger type; default would be 10pt

\usepackage[utf8]{inputenc} % set input encoding (not needed with XeLaTeX)

%%% Examples of Article customizations
% These packages are optional, depending whether you want the features they provide.
% See the LaTeX Companion or other references for full information.

%%% PAGE DIMENSIONS
\usepackage{geometry} % to change the page dimensions
\geometry{a4paper} % or letterpaper (US) or a5paper or....
\geometry{margin=1in} % for example, change the margins to 2 inches all round
% \geometry{landscape} % set up the page for landscape
%   read geometry.pdf for detailed page layout information

\usepackage{hyperref}
\usepackage{minted}
\usepackage{graphicx} % support the \includegraphics command and options

% \usepackage[parfill]{parskip} % Activate to begin paragraphs with an empty line rather than an indent

%%% PACKAGES
\usepackage{booktabs} % for much better looking tables
\usepackage{array} % for better arrays (eg matrices) in maths
\usepackage{paralist} % very flexible & customisable lists (eg. enumerate/itemize, etc.)
\usepackage{verbatim} % adds environment for commenting out blocks of text & for better verbatim
\usepackage{subfig} % make it possible to include more than one captioned figure/table in a single float
% These packages are all incorporated in the memoir class to one degree or another...

%%% HEADERS & FOOTERS
\usepackage{fancyhdr} % This should be set AFTER setting up the page geometry
\pagestyle{fancy} % options: empty , plain , fancy
\renewcommand{\headrulewidth}{0pt} % customise the layout...
\lhead{}\chead{}\rhead{}
\lfoot{}\cfoot{\thepage}\rfoot{}

%%% SECTION TITLE APPEARANCE
\usepackage{sectsty}
\allsectionsfont{\sffamily\mdseries\upshape} % (See the fntguide.pdf for font help)
% (This matches ConTeXt defaults)

%%% ToC (table of contents) APPEARANCE
\usepackage[nottoc,notlof,notlot]{tocbibind} % Put the bibliography in the ToC
\usepackage[titles,subfigure]{tocloft} % Alter the style of the Table of Contents
\renewcommand{\cftsecfont}{\rmfamily\mdseries\upshape}
\renewcommand{\cftsecpagefont}{\rmfamily\mdseries\upshape} % No bold!

\usepackage{color}

%%% END Article customizations

%%% The "real" document content comes below...

\title{Guide to using HINT}
\author{Yasuhiro SUZUKI}
\date{\today} % Activate to display a given date or no date (if empty),
         % otherwise the current date is printed 

\begin{document}
\maketitle

\section{Obtain the source code}
The source code of \texttt{HINT} and its relevant codes for pre-process and post-process are stored in GITHUB repository. In order to access the repository, you have to send the request to the author of the HINT code, Professor Yasuhiro Suzuki [\href{mailto:suzukiy@hiroshima-u.ac.jp}{suzukiy@hiroshima-u.ac.jp}].

\vspace{5mm}
You can download the source code via git clone;
\begin{minted}[linenos]{c}
$ git clone ssh://github.com/yasuhiro-suzuki/HINT3D.git
\end{minted}
\vspace{5mm}

In this tutorial, the \texttt{current}-branch of \texttt{HINT} code in the Github repository is used. 
The data used in the tutorial is stored in the server: \texttt{https://datashare.mpcdf.mpg.de/s/3iZlEPTRvFc8LPp}.

\section{Setting environment}
First of all, the environment of performing the \texttt{HINT} code has to be setted. 
The modules, especially \texttt{intel}, \texttt{intelmpi}, \texttt{hdf5}, \texttt{netcdf}, have to be loaded; 
\vspace{3mm}
\begin{minted}[linenos]{c}
$ module load intel/pe-xe-2020--binary
$ module load intelmpi/2020--binary
$ module load gnu/8.3.0
$ module load zlib/1.2.11--gnu--8.3.0
$ module load szip/2.1.1--gnu--8.3.0
$ module load hdf5
$ module load netcdf
$ module load netcdff

$ module list
Currently Loaded Modulefiles:
 1) profile/base                          
 2) env-skl/1.0                            
 3) intel/pe-xe-2020--binary               
 4) intelmpi/2020--binary                  
 5) gnu/8.3.0                              
 6) zlib/1.2.11--gnu--8.3.0                
 7) szip/2.1.1--gnu--8.3.0                 
 8) hdf5/1.12.2--intelmpi--2020--binary    
 9) netcdf/4.9.0--intelmpi--2020--binary   
10) netcdff/4.6.0--intelmpi--2020--binary  
\end{minted}
\vspace{3mm}


\section{Directory structure}

In this tutorial, we define \texttt{HINT} as our main project directory. 
Obtain the input data which contains the example of LHD from\texttt{https://datashare.mpcdf.mpg.de/s/3iZlEPTRvFc8LPp}, then place it in the main project directory, naming \texttt{DATA}.  
In the main directory, make directories of \texttt{build/} where you compile the codes. 
Sub-directories in the main directory is organized as follows:
\begin{verbatim}
$ DATA
$ bin
$ HINT3D/README.md
$ HINT3D/GPRTS
$ HINT3D/HINT
$ HINT3D/HMAG
$ HINT3D/MKFLX
$ HINT3D/MKLIM
$ HINT3D/MKVAC
$ build/GPRTS
$ build/HINT
$ build/HMAG
$ build/MKFLX
$ build/MKLIM
$ build/MKVAC

\end{verbatim}
The \texttt{MKVAC}, \texttt{MKFLX}, and \texttt{MKLIM} are preprocessing codes. \texttt{HMAG} and \texttt{GPRTS} are post-processing codes. These codes are called by shell scripts kept in \texttt{DATA} directory.

\section{Making input files for HINT}

\subsection{Introduction}
Making input files is relatively easy once you have some outputs from \texttt{VMEC}. The scripts are located in the \texttt{DATA} directory.

For pre-processing codes such as \texttt{MKVAC}, \texttt{MKFLX}, and \texttt{MKLIM}, the general procedure of compiling and execution is following:
\begin{description}
 \item[Step 1: VMEC input and output, and other inputs]\mbox{}\\
 Input files required are:
 \begin{itemize}
  \item \verb|mgrid.LHD_171x171x36| - this is a VMEC input file, it contains tables of the vacuum magnetic field of external coils. The \texttt{mgrid}-file is used by \texttt{MKVAC} to calculate the magnetic field in vacuum.  
  
  \item wout file - this is a VMEC output file, it contains information for making the normalized toroidal flux to prescribe the plasma pressure and toroidal current density. The \texttt{wout}-file is used by \texttt{MKFLX} to calculate the toroidal flux profile and the initial pressure profile is decided. 

  \item \verb|LHD_vessel_3D.dat| - this is an original data format file which contains the information of coordinates of the first-wall. The \texttt{vessel}-file is used by \texttt{MKLIM}.
 \end{itemize}
%THEN place all files in \texttt{MKMAG/DATA} directory.
The files of \texttt{mgrid}, \texttt{wout} and \texttt{vessel} are placed in \texttt{DATA} directory. 

 \item[Step 2: Compiling the input code]\mbox{}\\
  Make sure the source code to compile navigates to \texttt{build} directories and run \texttt{make}.

 \item[Step 3: Running the code]\mbox{}\\
  To run the code you need to navigate to \texttt{MAGFILE}, and then edit the shell script with the paths and names of the input files. Then execute the shell script.

 \item[Step 4: Making the controlling parameter file]\mbox{}\\
 To run the HINT code you need to specify the controlling parameters in the standard input.
 
\end{description}

\subsection{HINT code: main code for construction of 3D equilibrium}
First, compiling of the main code, \texttt{HINT}, is carried out. 

The compile should be done in \texttt{build} directory, by
\vspace{2mm}
\begin{minted}[linenos]{bash}
$ mkdir build
$ cd build
$ mkdir HINT
$ ln -s ../../HINT3D/HINT/*.f90 . 
$ ln -s ../../HINT3D/HINT/modules_stepa/fline_method_mod.f90 stepa_mod.f90
$ cp ../../HINT3D/HINT/Makefile.template Makefile
$ edit Makefile
$ make
\end{minted}
\vspace{2mm}

In the process of \texttt{edit Makefile} which in the 5-th line, the \texttt{Makefile} has to be editted according to your computer environment. 
For example, \texttt{Makefile} should be editted as below for executing on Marconi machine;
\vspace{2mm}
\begin{minted}[linenos]{make}
SHELL=/bin/sh
FC=mpiifort
LINK=mpiifort
FFLAGS=  -O2 -cpp -DNETCDF   ##-ffree-line-length-none -cpp
LDFLAGS=
MOD= -I$(HDF5_HOME)/include -I$(NETCDFF_HOME)/include
LIBS=-L$(HDF5_HOME)/lib -lhdf5hl_fortran -lhdf5_hl -lhdf5_fortran -lhdf5 -L$(NETCDF_HOME)/lib -L$(NETCDFF_HOME)/lib -lnetcdff -lnetcdf -L$(SZIP_LIB) -lsz -lz
PROGRAM=hint.exe
\end{minted}
\vspace{2mm}

When the compiling is successfully terminated, the directory of \texttt{build/HINT} has the execition file \texttt{hint.exe}. 

\subsection{MKVAC code: making the vacuum magnetic field}
The HINT needs the vacuum magnetic field as inputs. 
To make the input files follow these steps:
\begin{description}
 \item[Step 1: Compiling the code]\mbox{}\\
%  Run \texttt{make} in \texttt{build} directory.
Make a directory named \texttt{MKVAC} under the \texttt{build} directory, i.e. \texttt{build/MKVAC}. 

In the same way with the compile of HINT code, the source code should be linked as follows in the \texttt{build/MKVAC};
\vspace{2mm}
\begin{minted}[linenos]{bash}
$ ln -s ../../HINT3D/MKVAC/*.f90 .
$ cp ../../HINT3D/MKVAC/Makefile.gfortran Makefile
$ edit Makefile
$ make
\end{minted}
\vspace{2mm}

In the process of \texttt{edit Makefile} which in the 3rd line, the \texttt{Makefile} has to be editted according to your computer environment as like you did for compile of HINT code. 
For example, \texttt{Makefile} should be editted as below for executing on Marconi machine;
\vspace{2mm}
\begin{minted}[linenos]{make}
SHELL=/bin/sh
FC=ifort
LINK=ifort
FFLAGS= -convert big_endian -O2 -cpp -DNETCDF ##-ffree-line-length-none -cpp -DNETCDF
LDFLAGS=
LIBS=-L$(HDF5_HOME)/lib -lhdf5hl_fortran -lhdf5_hl -lhdf5_fortran -lhdf5 -L$(NETCDF_HOME)/lib -L$(NETCDFF_HOME)/lib -lnetcdff -lnetcdf -L$(SZIP_LIB) -lsz -lz
MOD=-I$(HDF5_HOME)/include -I$(NETCDFF_HOME)/include
PROGRAM=mkvac.exe
\end{minted}
\vspace{2mm}

Here in the 4th line, \texttt{FFLAGS=}, only LHD example case needs \texttt{FFLAGS= -convert big\_endian}. This is because the \texttt{mgrid}-file of LHD example case is fortran binary. 

When the compiling is successfully terminated, the directory of \texttt{build/MKVAC} has the execition file \texttt{mkvac.exe}. 


 \item[Step 2: Editing the shell script]\mbox{}\\
  You need to edit the shell script, \texttt{mkvac.csh}. Note set the following variables as:
   \begin{itemize}
    \item vac\_form: specify output file format;
     \begin{enumerate}
      \item HDF5 or hdf5: HDF5 format
      \item NETCDF or netcdf: netCDF4 format
      \item BIN or bin: Fortran binary format
     \end{enumerate}
    \item vac\_file: output file's name
    \item mtor: toroidal field period
    \item nr: grid number along $R$ \textcolor{red}{(must be $2^N$)}
    \item nz: grid number along $Z$ \textcolor{red}{(must be $2^N$)}
    \item ntor: grid number along $\phi$ \textcolor{red}{(must be $2^N$)}
    \item rmin: Minimum R in meter of rectangular computational box
    \item rmax: Maximum R in meter of rectangular computational box
    \item zmin: Minimum Z in meter of rectangular computational box
    \item zmax: Maximum Z in meter of rectangular computational box
    \item lsymmetry: logical flag to superimpose the stellarator symmetry \textcolor{red}{For stellarator, this flag should be set to T}
    \item mag\_form: specify the file format to read the code \textcolor{red}{For our case, `mgrid' must be used}
    \item extcur: external coil currents are specified
   \end{itemize}

 \item[Step 3: Running the code]\mbox{}\\
  Execute the shell script.
\vspace{-2mm}
\begin{minted}[linenos]{c}
$ ./mkvac.csh
\end{minted}
\vspace{-2mm}
\end{description}

The successful execution of \texttt{mkvac.csh} produces the output file as you have defined as \texttt{vac\_file}, for example, \texttt{r360\_g12538\_bq+100\_3.4x3.4\_256x256x64.vac.h5}. 

\subsection{MKFLX code: making the normalized toroidal flux distribution}
The HINT needs the normalized toroidal flux distribution as inputs. 
To make the input files follow these steps:
\begin{description}
 \item[Step 1: Compiling the code]\mbox{}\\
%%  Run \texttt{make} in \texttt{build} directory.
  Make a directory named \texttt{MKFLX} under the \texttt{build} directory, i.e. \texttt{build/MKFLX}. 

In the same way with the compile of HINT code, the source code should be linked as follows in the \texttt{build/MKFLX};
\vspace{2mm}
\begin{minted}[linenos]{bash}
$ ln -s ../../HINT3D/MKFLX/*.f90 .
$ cp ../../HINT3D/MKFLX/Makefile.gfortran Makefile
$ edit Makefile
$ make
\end{minted}
\vspace{2mm}

In the process of \texttt{edit Makefile} which in the 3rd line, the \texttt{Makefile} has to be edited according to your computer environment as like you did for compile of HINT code. 
For example, \texttt{Makefile} should be edited as below for executing on Marconi machine;
\vspace{2mm}
\begin{minted}[linenos]{make}
SHELL=/bin/sh
FC=ifort
LINK=ifort
FFLAGS= -convert big_endian -O2 -cpp -DNETCDF ##-ffree-line-length-none -cpp -DNETCDF -fbacktrace
LDFLAGS=-fbacktrace
LIBS=-L$(HDF5_HOME)/lib -lhdf5hl_fortran -lhdf5_hl -lhdf5_fortran -lhdf5 -L$(NETCDF_HOME)/lib -L$(NETCDFF_HOME)/lib -lnetcdff -lnetcdf -L$(SZIP_LIB) -lsz -lz
INC=
MOD=-I$(HDF5_HOME)/include -I$(NETCDFF_HOME)/include
PROGRAM=mkflx.exe
\end{minted}
\vspace{2mm}

When the compiling is successfully terminated, the directory of \texttt{build/MKFLX} has the execition file \texttt{mkflx.exe}. 

 \item[Step 2: Editing the shell script]\mbox{}\\
  You need to edit the shell script, \texttt{mkflx.csh}. Note set the following variables as:
   \begin{itemize}
    \item flx\_form: specify output file format;
     \begin{enumerate}
      \item HDF5 or hdf5: HDF5 format
      \item NETCDF or netcdf: netCDF4 format
      \item BIN or bin: Fortran binary format
     \end{enumerate}
    \item flx\_file: output file's name
    \item file\_format: specify wout file format
     \begin{enumerate}
      \item nc: netCDF3 or 4 format
      \item txt: ASCII format
     \end{enumerate}
    \item wout\_file: wout\_file's name 
    \item mtor: toroidal field period
    \item nr: grid number along $R$ \textcolor{red}{(must be $2^N$)}
    \item nz: grid number along $Z$ \textcolor{red}{(must be $2^N$)}
    \item ntor: grid number along $\phi$ \textcolor{red}{(must be $2^N$)}
    \item rmin: Minimum R in meter of rectangular computational box
    \item rmax: Maximum R in meter of rectangular computational box
    \item zmin: Minimum Z in meter of rectangular computational box
    \item zmax: Maximum Z in meter of rectangular computational box
   \end{itemize}

 \item[Step 3: Running the code]\mbox{}\\
  Execute the shell script.
\vspace{-2mm}
\begin{minted}[linenos]{bash}
$ ./mkflx.csh
\end{minted}
\vspace{-2mm}
\end{description}

The successful execution of \texttt{mkflx.csh} produces the output file as you have defined as \texttt{flx\_file}, for example, \texttt{r360\_g12538\_bq+100\_3.4x3.4\_256x256x64.flx.h5}. 



\subsection{MKLIM code: making limiter file}
The HINT needs the limiter file to decide the plasma boundary. The \texttt{MKLIM} code makes a mapping file in the inside and outside of the first wall. To make the limiter file follow these steps:
\begin{description}
 \item[Step 1: Compiling the code]\mbox{}\\
%  Run \texttt{make} in \texttt{build} directory.
  Make a directory named \texttt{MKLIM} under the \texttt{build} directory, i.e. \texttt{build/MKLIM}. 

In the same way with the compile of HINT code, the source code should be linked as follows in the \texttt{build/MKLIM};
\vspace{2mm}
\begin{minted}[linenos]{bash}
$ ln -s ../../HINT3D/MKLIM/*.f90 .
$ cp ../../HINT3D/MKLIM/Makefile.gfortran Makefile
$ edit Makefile
$ make
\end{minted}
\vspace{2mm}

In the process of \texttt{edit Makefile} which in the 3rd line, the \texttt{Makefile} has to be edited according to your computer environment as like you did for compile of HINT code. 
For example, \texttt{Makefile} should be edited as below for executing on Marconi machine;
\vspace{2mm}
\begin{minted}[linenos]{make}
SHELL=/bin/sh
FC=ifort
LINK=ifort
FFLAGS=-convert big_endian -O2 -cpp -DNETCDF ##-ffree-line-length-none -cpp -DNETCDF -fbacktrace
LDFLAGS=-fbacktrace
LIBS=-L$(HDF5_HOME)/lib -lhdf5hl_fortran -lhdf5_hl -lhdf5_fortran -lhdf5 -L$(NETCDF_HOME)/lib -L$(NETCDFF_HOME)/lib -lnetcdff -lnetcdf -L$(SZIP_LIB) -lsz -lz
INC=
MOD=-I$(HDF5_HOME)/include -I$(NETCDFF_HOME)/include
PROGRAM=mklim.exe
\end{minted}
\vspace{2mm}

When the compiling is successfully terminated, the directory of \texttt{build/MKLIM} has the execition file \texttt{mklim.exe}.

  
 \item[Step 2: Editing the shell script]\mbox{}\\
  You need to edit the shell script, \texttt{mklim.csh}. Note set the following variables as:
   \begin{itemize}
    \item lim\_form: specify output file format;
     \begin{enumerate}
      \item HDF5 or hdf5: HDF5 format
      \item NETCDF or netcdf: netCDF4 format
      \item BIN or bin: Fortran binary format
     \end{enumerate}
    \item lim\_file: output file's name
    \item vessel\_model: specify wout file format
     \begin{enumerate}
      \item 2D: Axisymmetric wall
      \item 3D: Non-axisymmetric wall
     \end{enumerate}
    \item vessel\_file: vessel file's name
    \item nr: grid number along $R$ \textcolor{red}{(Set the same value to \texttt{MKMAG})}
    \item nz: grid number along $Z$ \textcolor{red}{(Set the same value to \texttt{MKMAG})}
    \item ntor: grid number along $\phi$ \textcolor{red}{(Set the same value to \texttt{MKMAG})}
    \item mtor: toroidal field period \textcolor{red}{(Set the same value to \texttt{MKMAG})}
    \item rminb: Minimum R in meter of rectangular computational box \textcolor{red}{(must be same to \texttt{MKVAC} and \texttt{MKFLX})}
    \item rmaxb: Maximum R in meter of rectangular computational box \textcolor{red}{((must be same to \texttt{MKVAC} and \texttt{MKFLX})}
    \item zminb: Minimum Z in meter of rectangular computational box \textcolor{red}{(must be same to \texttt{MKVAC} and \texttt{MKFLX})}
    \item zmaxb: Maximum Z in meter of rectangular computational box \textcolor{red}{(must be same to \texttt{MKVAC} and \texttt{MKFLX})}
   \end{itemize}

 \item[Step 3: Running the code]\mbox{}\\
  Execute the shell script.
\vspace{-2mm}
\begin{minted}[linenos]{bash}
$ ./mklim.csh
\end{minted}
\vspace{-2mm}
\end{description}

The successful execution of \texttt{mklim.csh} produces the output file as you have defined as \texttt{lim\_file}, for example, \texttt{LHD\_vessel\_3.4x3.4\_256x256x64.limiter.h5}. 
  

\section{Running the HINT code}
\subsection{Making controlling parameter file}  

An example \verb|r360_g12538_bq+100_lc10_ppkf2.b020.input| in \texttt{DATA} is shown as follows:
\begin{verbatim}
!-----------------------------------------
! set MPI process numbers along R, Z, phi
!-----------------------------------------
&mpi_inp1
 iprocs = 1, ! along R idrection
 jprocs = 1, ! along Z direction
 kprocs = 4, ! along phi direction
&end
!--------------------------
! set mode to run the code
!--------------------------
&nlinp1
 lfile = .true.,
!-----------------------------------------------------
! "initial": clauclation set up from the vacuum field
! "follow" : following up job
!-----------------------------------------------------
 run_mode = 'initial',
!------------------------------------------------------------
! "file": normalized toroidal flux, s, is read from the file
! "calc": s is claulated from the pressure distribution
!------------------------------------------------------------
 flx_type = 'file',
!------------------------------------
! lresetp: redistribute the pressure
!------------------------------------
 lresetp  = .true.,
!----------------------------------------------------------------------
! "power series": specify the profile shape by 20th order polynominals
! "parabolic"   : specify the profile shape by parabolic functions
! "linear"      : linear interpolation from profile data
! "quadratic"   : quadatic interpolation (NOT implemented)
! "akima"       : Akima's cubic spline (NOT implemented)
!----------------------------------------------------------------------
 p_type   = 'parabolic',
&end
!------------------------
! set control parameters
!------------------------
&nlinp2
!--------------------------------------------------------------------
! beta:  toroidal beta value on axis
! bt0:   toroidal magnetic field to prescribe the beta
! am0:   polynominal coefficients to prescribe the pressure profile
! nstep: step number of overall loops
! npchg: step number to rump up the beta
! nsvae: interval to save the magnetic field and others
! npset: interval to redistribute the pressure
!--------------------------------------------------------------------
 beta  =  0.02,
 bt0   =  2.58,
 am0    = 0.01,  1.00,  1.00,
 nstep =  80,
 npchg =  1,
 nsave =  5,
 npset =  20,
&end
!-----------------------
! set coordinate system
!-----------------------
&nlinp3
!----------------------------------------------
! nr:    grid number along R
! nz:    grid number along Z
! ntor:  grid number along phi
! rminb: minimum R of the computational domain
! rmaxb: maximum R of the computational domain
! zminb: minimum Z of the computational domain
! zmaxb: maximum Z of the computational domain
! mtor:  toroidal field period
!----------------------------------------------
 nr     =  256,
 nz     =  256,
 ntor   =  64,
 rminb  =  2.2,
 rmaxb  =  5.6,
 zminb  = -1.7,
 zmaxb  =  1.7,
 mtor   =  10,
&end
&nlinp4
&end
&nlinp5
&end
!----------------------------------
! set control parameters in step-A
!----------------------------------
&stepa_inp1
!------------------------------------------------------------
! h_in:  step size of the Runge-Kutta for field line tracing
! lc_in: length of field line tracing
!------------------------------------------------------------
 h_in  =  1.0E-02,
 lc_in =  1.0E+01,
&end
!----------------------------------
! set control parameters in step-B
!----------------------------------
&stepb_inp1
!---------------------------------------------------------
! lhistoryb: switch to write debugging information
! nstepb:    step number to evole MHD equations
! nrefb:     interval to save refrences for smoothing
! nsoothb:   interval to apply smoothing filter
! nenergyb:  interval to caluclate debugging information
! dt_b:      time step to evole MHD equations
! eta0:      resistivity in the plasma
! nu1:       viscosity
!---------------------------------------------------------
 lhistoryb = .true.,
 nstepb    = 1000,
 nrefb     = 1000,
 nsmoothb  = 1,
 nenergyb  = 100,
 dt_b      = 1.0E-03,
 eta0      = 1.0E-04,
 nu1       = 1.0E-08
&end
&fopen_stepb
 historyb_file = 'r360_g12538_bq+100_lc10_ppkf2.b020.historyb',
&end
&fopen
 file_type     = 'hdf5',
 vac_file      = 'r360_g12538_bq+100_3.4x3.4_256x256x64.vac.h5',
 flx_file      = 'r360_g12538_bq+100_3.4x3.4_256x256x64.flx.h5',
 limiter_file  = 'LHD_vessel_3.4x3.4_256x256x64.limiter.h5',
 mag_file      = 'r360_g12538_bq+100_lc10_ppkf2.b020.mag.h5',
&end
\end{verbatim}
You must specify the following parameters as follows:
\begin{itemize}
 \item \texttt{iprocs} : MPI processes number of $R$-direction
 \item \texttt{jprocs} : MPI processes number of $Z$-direction
 \item \texttt{kprocs} : MPI processes number of $\phi$-direction
 \item \texttt{beta} : toroidal beta
 \item \texttt{bt0} : toroidal field to specify the toroidal beta
 \item \texttt{nstep} : total iteration number
 \item \texttt{npchg} : step number to rump up the beta
 \item \texttt{nsave} : interval to save the magnetic field and others
 \item \texttt{npset} : interval to redistribute the pressure
 \item \texttt{p\_type} : \texttt{HINT} specifies the pressure profile shape as a function of the normalized toroidal flux. The parameter \texttt{p\_type} set to the type of functions.
 \item \texttt{am0} : coefficients of the profile function
 \item \texttt{nr} : grid numbers of $R$-direction. \textcolor{red}{Set the same value to the input file by \texttt{MKVAC}}.
 \item \texttt{nz} : grid numbers of $Z$-direction. \textcolor{red}{Set the same value to the input file by \texttt{MKVAC}}.
 \item \texttt{ntor} : grid numbers of $\phi$-direction. \textcolor{red}{Set the same value to the input file by \texttt{MKVAC}}.
 \item \texttt{mtor} : toroidal field period. \textcolor{red}{Set the same value to the input file by \texttt{MKVAC}}.
 \item \texttt{h\_in} : step size of the Runge-Kutta for field line tracing.
 \item \texttt{lc\_in} : length of field line tracing.
 \item \texttt{eta0} : resistivity.
 \item \texttt{hisotryb\_file} : specify the file name on the scratch directory.
 \item \texttt{file\_type} : specify the input and output file formats.
  \begin{enumerate}
   \item  HDF5 or hdf5: HDF5 format
   \item NETCDF or netcdf: netCDF4 format
   \item BIN or bin: Fortran binary format
  \end{enumerate}
 \item \texttt{vac\_file} : specify the file name in the project directory.
 \item \texttt{flx\_file} : specify the file name in the project directory.
 \item \texttt{limiter\_file} : specify the file name in the project directory.
 \item \texttt{mag\_file} : specify the file name on the scratch directory.
\end{itemize}

\subsection{Running HINT}

\textcolor{red}{This part describes only the test job. This part will be updated.}

\begin{description}

 \item[Step 1: Preparing input files and shell script]\mbox{}\\
  You need to prepare the working directory. 
  Here, in this tutorial, the working directory is named as \texttt{run\_LHD}. 

  In the working directory, place the output files of \texttt{MKVAC}, \texttt{MKFLX} and \texttt{MKLIM}, i.e. \texttt{r360\_g12538\_bq+100\_3.4x3.4\_256x256x64.vac.h5}, \texttt{r360\_g12538\_bq+100\_3.4x3.4\_256x256x64.flx.h5}, and \texttt{LHD\_vessel\_3.4x3.4\_256x256x64.limiter.h5}, respectively.  

  Also, copy an input file, \texttt{r360\_g12538\_bq+100\_lc10\_ppkf2.b020.input} and shell script \texttt{go001.csh}, from \texttt{DATA} directory to the working directory. 
  
  %Please copy \texttt{/data/lng/suzuki/LHD/DATA/}. After making working directory you have to copy an input file, \texttt{r360\_g12538\_bq+100\_lc10\_ppkf2.b020.input}, and shell script, \texttt{go001.csh}. These are examples to run HINT. THEN edit these files to match your environment (i.e. file name, scratch directory, and so on...)

 \item[Step 2: Submitting batch job]\mbox{}\\

 Here, instead of \texttt{go001.csh}, the batch file for Marconi machine is introduced. 
 The batch file has to request the same number of MPI process which you defined in the input file, \texttt{r360\_g12538\_bq+100\_lc10\_ppkf2.b020.input}. Then, the mpi execution has to be launched reading the input file;
 \begin{verbatim}
 $  mpirun ../build/HINT/hint.exe < r360_g12538_bq+100_lc10_ppkf2.b020.input > out_hint
\end{verbatim}
  
  Submit a job as:
  \begin{verbatim}
   sbatch job_marconi.sh
  \end{verbatim}
%  \begin{verbatim}
%   pjsub go001.csh
%  \end{verbatim}
  If the cluster is not clouded, the calculation will be finished within a few minutes (depending on \texttt{nstep}).
\end{description}


The successfull execution of the HINT code produces the following files;
\texttt{r360\_g12538\_bq+100\_lc10\_ppkf2.b020.historyb} and \texttt{r360\_g12538\_bq+100\_lc10\_ppkf2.b020.mag.h5}.


\section{Analysis using HMAG}
\subsection{Introduction}
HMAG is to follow the field lines to analyze the magnetic surfaces. 

\subsection{Compile}
  Make a directory named \texttt{HMAG} under the \texttt{build} directory, i.e. \texttt{build/HMAG}. 

In the same way with the compile of HINT code, the source code should be linked as follows in the \texttt{build/HMAG};
\vspace{2mm}
\begin{minted}[linenos]{bash}
$ ln -s ../../HINT3D/HMAG/*.f90 .
$ cp ../../HINT3D/HMAG/Makefile.gfortran Makefile
$ edit Makefile
$ make
\end{minted}
\vspace{2mm}

In the process of \texttt{edit Makefile} which in the 3rd line, the \texttt{Makefile} has to be edited according to your computer environment as like you did for compile of HINT code. 
For example, \texttt{Makefile} should be edited as below for executing on Marconi machine;
\vspace{2mm}
\begin{minted}[linenos]{make}
SHELL=/bin/sh
FC=ifort
LINK=ifort
FFLAGS=-convert big_endian -O2 -cpp -DNETCDF  ##-ffree-line-length-none -cpp -DPLPLOT -DNETCDF
LDFLAGS=
MOD=-I$(HDF5_HOME)/include -I$(NETCDFF_HOME)/include
LIBS=-L$(HDF5_HOME)/lib -lhdf5hl_fortran -lhdf5_hl -lhdf5_fortran -lhdf5 -L$(NETCDF_HOME)/lib -L$(NETCDFF_HOME)/lib -lnetcdff -lnetcdf -L$(SZIP_LIB) -lsz -lz
PROGRAM=hmag.exe
\end{minted}
\vspace{2mm}

When the compiling is successfully terminated, the directory of \texttt{build/HMAG} has the execition file \texttt{hmag.exe}.

\subsection{Execution - calculation of $\beta$-field}

Copy the shell script \texttt{hmag.csh} from \texttt{DATA} directory to working directory. 
Modify the shell script \texttt{hmag.csh} according to the structure of directory. 

The HMAG code requires the information of the geometry of vacumm vessel, therefore, you need to copy \texttt{LHD\_vessel\_3D.dat} from \texttt{DATA};
\begin{minted}[linenos]{bash}
$ cp ../DATA/LHD_vessel_3D.dat .
\end{minted}

Execute the shell script.
\vspace{-2mm}
\begin{minted}[linenos]{bash}
$ ./hmag.csh
\end{minted}
\vspace{5mm}


The successful execution of \texttt{hmag.csh} produces the output files, for example, \texttt{r360\_g12538\_bq+100\_lc10\_ppkf2.b020.mag.h5}, \texttt{fort.60}, and \texttt{r360\_g12538\_bq+100\_lc10\_ppkf2.b020.1.plot}-\texttt{r360\_g12538\_bq+100\_lc10\_ppkf2.b020.8.plot}. 

Files \texttt{r360\_g12538\_bq+100\_lc10\_ppkf2.b020.1.plot} etc. contain the data of the poincare plot. 
In order to plot, the easiest way is to use \texttt{gnuplot};
\vspace{0mm}
\begin{minted}[linenos]{bash}
$ gnuplot
gnuplot> plot 'r360_g12538_bq+100_lc10_ppkf2.b020.1.plot'
\end{minted}
\vspace{5mm}

The file \texttt{fort.60} contains the information of the plasma minor radius $<r>$, rotational transform $\iota$. 

\subsection{Calculation of vacuum field}
The vacuum field (i.e. before execution of HINT) can be computed using \texttt{HMAG}. 

You need the files of the output files of \texttt{MKVAC}, \texttt{MKFLX} and \texttt{MKLIM}, i.e.
\begin{itemize}
    \item \texttt{r360\_g12538\_bq+100\_3.4x3.4\_256x256x64.vac.h5}, 
    \item \texttt{r360\_g12538\_bq+100\_3.4x3.4\_256x256x64.flx.h5},
    \item \texttt{LHD\_vessel\_3.4x3.4\_256x256x64.limiter.h5}
    \item \texttt{LHD\_vessel\_3D.dat}.
\end{itemize}

Set the parameters such as \texttt{rax0}, etc. Then, execute by;
\vspace{-2mm}
\begin{minted}[linenos]{bash}
$ ./vac.csh 
\end{minted}
\vspace{2mm}


\section{Preparation of MIPS input from HINT}
Output of HINT code can be converted into input of MIPS code. 
The code which converts output of HINT code into input of MIPS code is named \texttt{MAGVAL}, and it is stored in \texttt{DATA}. 

First, you have to set the configuration of Makefile;
For example, \texttt{Makefile} should be edited as below for executing on Marconi machine;
\vspace{2mm}
\begin{minted}[linenos]{make}
SHELL=/bin/sh
FC=ifort
LINK=ifort
FFLAGS=-byteswapio -mcmodel=medium -O2 -cpp -DNETCDF -traceback
LDFLAGS=-traceback
MOD=-I$(HDF5_HOME)/include -I$(NETCDFF_HOME)/include
LIBS=
LIBS=-L$(HDF5_HOME)/lib -lhdf5hl_fortran -lhdf5_hl -lhdf5_fortran -lhdf5 -L$(NETCDF_HOME)/lib -L$(NETCDFF_HOME)/lib -lnetcdff -lnetcdf -L$(SZIP_LIB) -lsz -lz
PROGRAM=magval.exe
\end{minted}
\vspace{2mm}

The successful compiling gives the output file \texttt{magval.exe}. 

Place the \texttt{magval.exe} and the shell script \texttt{go.csh} in the working directory of HINT calculation. 
Modify the shell script \texttt{go.csh} according to the computer environment. 

Note, the HINT calculation has been performed with \texttt{rmin=2.2},  \texttt{rmax=5.6},  \texttt{zmin=-1.7} and \texttt{zmin=1.7}. 
For MIPS calculation, the simulation domain has to extract the inner part of HINT calculation, because MIPS code wants to avoid the coil positions, \texttt{rmin=2.8},  \texttt{rmax=4.8},  \texttt{zmin=-1.0} and \texttt{zmin=1.0}.


Then, execute the shell script;
\vspace{0mm}
\begin{minted}[linenos]{bash}
$ ./go.csh 
\end{minted}
\vspace{5mm}

The successful execution of \texttt{magval.exe} gives the output files, 
\begin{itemize}
    \item \texttt{fort.400}
    \item \texttt{r360\_g12538\_bq+100\_lc10\_ppkf2.b020.mips\_eq}. 
\end{itemize}


\section{Appendix}
\subsection{Execution of TJ-II}
\subsection{MKVAC code: making the vacuum magnetic field}

The compiling of the \texttt{MKVAC} code has been already done in the tutorial of LHD case. 
\begin{description}
 \item[Step 2: Editing the shell script]\mbox{}\\
  You need to edit the shell script, \texttt{mkvac2.csh}. 

 \item[Step 3: Running the code]\mbox{}\\
  Execute the shell script.
\vspace{-2mm}
\begin{minted}[linenos]{bash}
$ ./mkvac2.csh
\end{minted}
\vspace{-2mm}
\end{description}

The successful execution of \texttt{mkvac2.csh} produces the output file as you have defined as \texttt{vac\_file}, for example, \texttt{TJ-II\_tutorial.vac.h5}. 



\subsection{MKFLX code: making the normalized toroidal flux distribution}
The compilation has been alredady done. 
\begin{description}


 \item[Step 2: Editing the shell script]\mbox{}\\
  You need to edit the shell script, \texttt{mkflx.csh}. 

 \item[Step 3: Running the code]\mbox{}\\
  Execute the shell script.
\vspace{-2mm}
\begin{minted}[linenos]{bash}
$ ./mkflx.csh
\end{minted}
\vspace{-2mm}
\end{description}

The successful execution of \texttt{mkflx.csh} produces the output file as you have defined as \texttt{flx\_file}, for example, \texttt{TJ-II\_tutorial.flx.h5}. 


\subsection{MKLIM code: making limiter file}
The compilation has been done. 
\begin{description}

  
 \item[Step 2: Editing the shell script]\mbox{}\\
  You need to edit the shell script, \texttt{mklim.csh}. 
  
 \item[Step 3: Running the code]\mbox{}\\
  Execute the shell script.
\vspace{-2mm}
\begin{minted}[linenos]{bash}
$ ./mklim.csh
\end{minted}
\vspace{-2mm}
\end{description}

The successful execution of \texttt{mklim.csh} produces the output file as you have defined as \texttt{lim\_file}, for example, \texttt{TJ-II\_tutorial.limiter.h5}. 
  


\subsection{Running HINT}


\begin{description}

 \item[Step 1: Preparing input files and shell script]\mbox{}\\
  You need to prepare the working directory. 
  Here, in this tutorial, the working directory is named as \texttt{run\_TJ-II}. 

  In the working directory, place the output files of \texttt{MKVAC}, \texttt{MKFLX} and \texttt{MKLIM}, i.e. \texttt{TJ-II\_tutorial.vac.h5}, \texttt{TJ-II\_tutorial.flx.h5}, and \texttt{TJ-II\_tutorial.limiter.h5}, respectively.  

  Also, copy an input file, \texttt{TJ-II.input} (here, the name has been changed from \texttt{r360\_g12538\_bq+100\_lc10\_ppkf2.b020.input}) and shell script \texttt{go001.csh}, from \texttt{DATA} directory to the working directory. 

 \item[Step 2: Submitting batch job]\mbox{}\\

 The batch file has to request the same number of MPI process which you defined in the input file, \texttt{TJ-II.input}. Then, the mpi execution has to be launched reading the input file;
 \begin{verbatim}
 $  mpirun ../build/HINT/hint.exe < TJ-II.input > out_hint
\end{verbatim}
  
  Submit a job as:
  \begin{verbatim}
   sbatch job_marconi.sh
  \end{verbatim}
%  \begin{verbatim}
%   pjsub go001.csh
%  \end{verbatim}
  If the cluster is not clouded, the calculation will be finished within a few minutes (depending on \texttt{nstep}).
\end{description}


The successful execution of the HINT code produces the following files;
\texttt{TJ-II\_tutorial.historyb} and \texttt{TJ-II\_tutorial.mag.h5}.


\subsection{Calculation of vacuum field}
The vacuum field (i.e. before execution of HINT) can be computed using \texttt{HMAG}. 

You need the files of the output files of \texttt{MKVAC}, \texttt{MKFLX} and \texttt{MKLIM}, i.e.
\begin{itemize}
    \item \texttt{TJ-II\_tutorial.vac.h5}, 
    \item \texttt{TJ-II\_tutorial.flx.h5},
    \item \texttt{TJ-II\_tutorial.limiter.h5}
    \item \texttt{vessel.dat}.
\end{itemize}

Set the parameters such as \texttt{rax0}, etc. Then, execute by;
\vspace{-2mm}
\begin{minted}[linenos]{bash}
$ ./vac.csh 
\end{minted}
\vspace{2mm}


\subsection{HMAG calculation of $\beta$-field}

Copy the shell script \texttt{hmag.csh} from \texttt{DATA} directory to working directory. 
Modify the shell script \texttt{hmag.csh} according to the structure of directory. 

The HMAG code requires the information of the geometry of vacumm vessel, therefore, you need to copy \texttt{LHD\_vessel\_3D.dat} from \texttt{DATA};
\begin{minted}[linenos]{bash}
$ cp ../DATA/LHD_vessel_3D.dat .
\end{minted}

Execute the shell script.
\vspace{-2mm}
\begin{minted}[linenos]{bash}
$ ./hmag.csh
\end{minted}
\vspace{5mm}


The successful execution of \texttt{hmag.csh} produces the output files, for example, \texttt{r360\_g12538\_bq+100\_lc10\_ppkf2.b020.mag.h5}, \texttt{fort.60}, and \texttt{r360\_g12538\_bq+100\_lc10\_ppkf2.b020.1.plot}-\texttt{r360\_g12538\_bq+100\_lc10\_ppkf2.b020.8.plot}. 

Files \texttt{r360\_g12538\_bq+100\_lc10\_ppkf2.b020.1.plot} etc. contain the data of the poincare plot. 
In order to plot, the easiest way is to use \texttt{gnuplot};
\vspace{0mm}
\begin{minted}[linenos]{bash}
$ gnuplot
gnuplot> plot 'r360_g12538_bq+100_lc10_ppkf2.b020.1.plot'
\end{minted}
\vspace{5mm}

The file \texttt{fort.60} contains the information of the plasma minor radius $<r>$, rotational transform $\iota$. 



\end{document}
